% !TeX program = pdflatex
% CoderUJB, Equipe 6: 32 slides, including cover, contents, six section pages and all 24 content slides.
% Derived from the supplied 24-slide French PPTX.
% Visual reference supplied by the user: sildeS2Groupe1.pdf.
% Its native Beamer settings are available in rapportS2_beamer.tex.
% Only layout settings are reused; the source project's content and identity are not imported.
\documentclass[aspectratio=169]{beamer}
\mode<presentation>
{
  \usetheme{default}
  \usecolortheme{default}
  \usefonttheme{default}
  \setbeamertemplate{navigation symbols}{}
  \setbeamertemplate{caption}[numbered]
  \setbeamertemplate{footline}{%
    \leavevmode%
    \hbox{%
      \begin{beamercolorbox}[wd=\paperwidth,ht=2.6ex,dp=1.2ex,rightskip=1.2em]{author in head/foot}%
        \hfill\usebeamerfont{author in head/foot}\insertframenumber
      \end{beamercolorbox}%
    }%
  }
}
\usepackage[french]{babel}
\usepackage[utf8]{inputenc}
\usepackage[T1]{fontenc}
\usepackage{lmodern}
\usepackage{amsmath}
\usepackage{amssymb}
\usepackage{booktabs}
\usepackage{array}
\usepackage{ragged2e}

% Give slide headings the same left margin as the body and extra top space.
\setbeamertemplate{frametitle}{%
  \nointerlineskip
  \vspace*{0.70cm}%
  \begin{beamercolorbox}[wd=\paperwidth,leftskip=1cm,rightskip=1cm,sep=0pt]{frametitle}
    \usebeamerfont{frametitle}\strut\insertframetitle\strut\par
  \end{beamercolorbox}%
  \vspace*{0.12cm}%
}

\AtBeginEnvironment{frame}{\justifying}
\AtBeginEnvironment{block}{\justifying}
\AtBeginEnvironment{alertblock}{\justifying}
\AtBeginEnvironment{exampleblock}{\justifying}
\addtobeamertemplate{itemize/enumerate body begin}{\justifying}{}
\addtobeamertemplate{itemize/enumerate subbody begin}{\justifying}{}
\addtobeamertemplate{itemize/enumerate subsubbody begin}{\justifying}{}
\setbeameroption{hide notes}
\hypersetup{
  pdftitle={CoderUJB -- Equipe 6},
  pdfauthor={Wei WEI, Hao QIAN, Jianye SHI, Xiang BIAN, Zijie Fang, Yupan Huang},
  pdfsubject={Presentation de l'article CoderUJB, ISSTA 2024},
  pdfkeywords={CoderUJB, Java, LLM, benchmark, ISSTA 2024}
}
% Presenter names and order supplied by the user on 2026-11-05.
\newcommand{\PresentationSpeakers}{\begin{tabular}{ccc}Wei WEI & Hao QIAN & Jianye SHI\\[0.4em]Xiang BIAN & Zijie Fang & Yupan Huang\end{tabular}}
\title[CoderUJB]{CoderUJB: An Executable and Unified Java Benchmark for Practical Programming Scenarios}
\subtitle{Présentation d'article -- ISSTA 2024}
\author{{\small Présenté par}\\[0.5em]\PresentationSpeakers}
\institute{Master CHPS -- UVSQ\\EDP\\[0.8em]{\scriptsize Article de Zhengran Zeng, Yidong Wang, Rui Xie, Wei Ye et Shikun Zhang}}
\date{}

% Six presentation sections, in the presenter order confirmed by the group.
\newcommand{\SpeakerForNumber}[1]{%
  \ifcase#1\relax
  \or Wei WEI%
  \or Hao QIAN%
  \or Jianye SHI%
  \or Xiang BIAN%
  \or Zijie Fang%
  \or Yupan Huang%
  \fi
}
\setbeamertemplate{section in toc}{%
  \leavevmode\usebeamercolor[fg]{structure}%
  \inserttocsectionnumber.\ \inserttocsection\par
}
\AtBeginSection[]{%
  \begin{frame}
    \centering
    {\usebeamercolor[fg]{structure}\Large
      \insertsectionnumber.\ \insertsectionhead\par}
    \vspace{1.3em}
    {\large\SpeakerForNumber{\value{section}}\par}
  \end{frame}
}
\begin{document}




% Slide 01 / 32; intervenant 1

\begin{frame}
\titlepage
\note{
INTERVENANT 1 — Wei WEI — Contexte de la recherche

DIAPOSITIVE 1 / 32
}
\end{frame}

\begin{frame}{Sommaire}
  \tableofcontents[hideallsubsections]
\end{frame}

\section{Contexte de la recherche}

% Slide 04 / 32; intervenant 1
\begin{frame}{Comment évaluer une IA qui code ?}
\small
\raggedright
{\large Un code qui semble correct\\[0.3em]
peut échouer à l’exécution.\par}
\vspace{1.4em}
{\large\structure{\textbf{CODE + TESTS}}\par}
\vspace{1.4em}
Benchmark : des tâches et des règles communes

\note{
INTERVENANT 1 — Wei WEI — Contexte de la recherche

DIAPOSITIVE 4 / 32

Bonjour à tous. Aujourd’hui, nous allons présenter l’article CoderUJB, publié en 2024. Je vais d’abord expliquer le contexte de cette étude.

Les grands modèles de langage, qu’on appelle aussi les LLM, peuvent aider à programmer. Ils peuvent, par exemple, proposer du code ou aider à corriger une erreur.

Mais comment savoir si leurs réponses sont bonnes ? [Pause]

Un code peut sembler correct quand on le lit. Pourtant, il peut échouer quand on l’exécute. Il faut donc le tester.

Pour comparer les modèles, les chercheurs utilisent des benchmarks. Un benchmark, c’est un ensemble de tâches avec des règles communes pour évaluer les réponses.

REPÈRES : IA  /  code  /  tests  /  benchmark

SOURCE : Zeng et al. (2024), CoderUJB: An Executable and Unified Java Benchmark for Practical Programming Scenarios. Introduction, tableau 1 et §2, pp. 124-126. DOI: 10.1145/3650212.3652115.

Les chiffres reprennent l'article et n'ont pas été reproduits par notre groupe. Partie 1 : 407 mots ; cible inférieure à 4 minutes, à vérifier par répétition chronométrée.
}
\end{frame}



% Slide 05 / 32; intervenant 1

\begin{frame}{Un exercice isolé a ses limites}
\small
\begin{columns}[T,onlytextwidth]
\begin{column}{0.46\textwidth}
{\usebeamercolor[fg]{structure}\normalsize\textbf{HUMANEVAL}\par}
\medskip
164 exercices Python\par
\vspace{0.8em}
\begin{itemize}
  \item Une fonction isolée
  \item Vérification par des tests
\end{itemize}
\end{column}
\begin{column}{0.48\textwidth}
{\usebeamercolor[fg]{structure}\normalsize\textbf{DANS UN VRAI PROJET}\par}
\medskip
Du code déjà présent\par
\vspace{0.8em}
\begin{itemize}
  \item D’autres fichiers
  \item Des bibliothèques à utiliser
\end{itemize}
\end{column}
\end{columns}

\note{
INTERVENANT 1 — Wei WEI — Contexte de la recherche

DIAPOSITIVE 5 / 32

Un exemple connu est HumanEval. Il contient cent soixante-quatre exercices en Python. Le modèle doit compléter une fonction à partir d’une description. Ensuite, des tests vérifient le résultat.

C’est utile, car on exécute vraiment le code. [Pause]

Mais les exercices restent assez isolés. Dans un vrai projet, une fonction peut utiliser d’autres fichiers ou des bibliothèques. Le modèle doit donc tenir compte du code qui existe déjà.

On peut comparer cela à un exercice de cours et à un travail dans un logiciel complet. Réussir le premier ne garantit pas qu’on réussira le second. C’est une première difficulté pour l’évaluation.

REPÈRES : HumanEval  /  fonction  /  projet  /  dépendances

SOURCE : Zeng et al. (2024), CoderUJB: An Executable and Unified Java Benchmark for Practical Programming Scenarios. Introduction, tableau 1 et §2, pp. 124-126. DOI: 10.1145/3650212.3652115.

Les chiffres reprennent l'article et n'ont pas été reproduits par notre groupe. Partie 1 : 407 mots ; cible inférieure à 4 minutes, à vérifier par répétition chronométrée.
}
\end{frame}

% Slide 06 / 32; intervenant 1

\begin{frame}{Développer, c’est faire plusieurs tâches}
\small
\begin{columns}[T,onlytextwidth]
\begin{column}{0.22\textwidth}\structure{\large\textbf{Écrire}}\end{column}
\begin{column}{0.22\textwidth}\structure{\large\textbf{Tester}}\end{column}
\begin{column}{0.22\textwidth}\structure{\large\textbf{Détecter}}\end{column}
\begin{column}{0.22\textwidth}\structure{\large\textbf{Corriger}}\end{column}
\end{columns}
\vspace{1.8em}
Une seule tâche donne une vision limitée.\par
\vspace{1.6em}
CodeXGLUE : ressemblance entre les textes\par
\smallskip
XCodeEval : problèmes de concours\par

\note{
INTERVENANT 1 — Wei WEI — Contexte de la recherche

DIAPOSITIVE 6 / 32

Il y a aussi une deuxième difficulté. Le travail d’un développeur comprend plusieurs types de tâches. Il faut écrire du code, mais aussi créer des tests, trouver des erreurs et les corriger.

Un modèle peut être bon pour une tâche et moins bon pour une autre. Une seule tâche donne donc une vision limitée de ses capacités. [Pause]

Certains benchmarks proposent plusieurs tâches. Mais, selon les auteurs, ils ont d’autres limites. Par exemple, CodeXGLUE utilise des mesures de ressemblance entre les textes. XCodeEval utilise des problèmes de concours.

Ces approches ne réunissent donc pas toutes les conditions recherchées pour évaluer le travail dans un vrai projet.

REPÈRES : écrire  /  tester  /  détecter  /  corriger

SOURCE : Zeng et al. (2024), CoderUJB: An Executable and Unified Java Benchmark for Practical Programming Scenarios. Introduction, tableau 1 et §2, pp. 124-126. DOI: 10.1145/3650212.3652115.

Les chiffres reprennent l'article et n'ont pas été reproduits par notre groupe. Partie 1 : 407 mots ; cible inférieure à 4 minutes, à vérifier par répétition chronométrée.
}
\end{frame}

% Slide 07 / 32; intervenant 1

\begin{frame}{Le besoin : une évaluation plus réaliste}
\small
{\renewcommand{\arraystretch}{2.1}
\begin{tabular}{@{}l@{\hspace{1.1em}}l@{\hspace{1.8em}}l@{}}
\structure{01} & \structure{\textbf{Exécuter}} & Vérifier le fonctionnement \\
\structure{02} & \structure{\textbf{Vrais projets}} & Code et environnement \\
\structure{03} & \structure{\textbf{Plusieurs tâches}} & Des capacités variées
\end{tabular}\par}
\vspace{1.3em}
CoderUJB : un benchmark Java pour des situations réelles

\note{
INTERVENANT 1 — Wei WEI — Contexte de la recherche

DIAPOSITIVE 7 / 32

Les auteurs cherchent donc à réunir trois éléments dans une même évaluation.

D’abord, exécuter les solutions pour vérifier leur fonctionnement. Ensuite, utiliser de vrais projets, avec leur code et leur environnement. Enfin, proposer plusieurs types de tâches. [Pause]

C’est dans ce contexte qu’ils présentent CoderUJB. Ils choisissent Java, un langage très utilisé dans le développement de logiciels.

La question de départ est donc simple : comment mieux évaluer les capacités des modèles dans des situations proches du travail réel ?

Voilà la motivation de cette étude. Je laisse maintenant la parole à la personne suivante, qui va présenter les contributions de l’article.

REPÈRES : exécution  /  projets  /  tâches  /  Java

SOURCE : Zeng et al. (2024), CoderUJB: An Executable and Unified Java Benchmark for Practical Programming Scenarios. Introduction, tableau 1 et §2, pp. 124-126. DOI: 10.1145/3650212.3652115.

Les chiffres reprennent l'article et n'ont pas été reproduits par notre groupe. Partie 1 : 407 mots ; cible inférieure à 4 minutes, à vérifier par répétition chronométrée.
}
\end{frame}

\section{Contributions et méthode}

% Slide 09 / 32; intervenant 2

\begin{frame}{Un benchmark commun pour cinq tâches}
\small
\begin{center}
\renewcommand{\arraystretch}{1.5}
\begin{tabular}{ccc}
{\Large\structure{2\,239}} & {\Large\structure{17}} & {\Large\structure{5}} \\
questions & projets Java & tâches
\end{tabular}
\end{center}
\medskip
{\small Source commune : Defects4J. Un benchmark et une étude comparative.\par}

\note{
INTERVENANT 2 — Hao QIAN — Contributions et méthode

DIAPOSITIVE 9 / 32

Merci. Je vais présenter ce que les auteurs proposent et comment ils évaluent les modèles.

CoderUJB contient deux mille deux cent trente-neuf questions. Elles viennent de dix-sept projets Java disponibles dans Defects4J. Cette source contient des programmes, des bugs et des versions corrigées.

Le premier apport est de réunir cinq tâches dans un même cadre. Le deuxième est de fournir le projet et son environnement d’exécution. Le troisième est une étude comparative de plusieurs modèles. 

REPÈRES : 2 239 questions / 17 projets / Defects4J

SOURCE : Zeng et al. (2024), CoderUJB: An Executable and Unified Java Benchmark for Practical Programming Scenarios. Résumé, §3.2, pp. 124, 127. DOI: 10.1145/3650212.3652115.

Les chiffres reprennent l'article et n'ont pas été reproduits par notre groupe. Partie 2 : 337 mots ; cible inférieure à 4 minutes, à vérifier par répétition chronométrée.
}
\end{frame}

% Slide 10 / 32; intervenant 2

\begin{frame}{Les cinq tâches de CoderUJB}
\small
\begin{center}
\small
\renewcommand{\arraystretch}{1.35}
\begin{tabular}{@{}lp{6.1cm}r@{}}
\toprule
\textbf{Tâche} & \textbf{Action demandée} & \textbf{Questions} \\
\midrule
FCG & Compléter une fonction & 238 \\
CTG & Créer un test à partir du code & 140 \\
ITG & Reproduire un bug signalé & 451 \\
APR & Corriger une fonction & 470 \\
DD & Détecter un bug : oui / non & 940 \\
\bottomrule
\end{tabular}
\end{center}
\medskip
{\small Total : 2 239 questions. DD comprend 470 versions défectueuses et 470 versions corrigées.\par}

\note{
INTERVENANT 2 — Hao QIAN — Contributions et méthode

DIAPOSITIVE 10 / 32

Voici les cinq tâches. La première consiste à compléter une fonction à partir de sa description. La deuxième demande d’écrire un test à partir du code. La troisième part d’un signalement de bug : le test doit reproduire ce problème.

La quatrième consiste à corriger une fonction défectueuse. Enfin, la cinquième demande de dire si une fonction contient un bug.

Les deux tâches de test ont donc des entrées différentes : le code dans un cas, le signalement du problème dans l’autre. 

REPÈRES : FCG / CTG / ITG / APR / DD

SOURCE : Zeng et al. (2024), CoderUJB: An Executable and Unified Java Benchmark for Practical Programming Scenarios. §3.1-3.3, pp. 126-128. DOI: 10.1145/3650212.3652115.

Les chiffres reprennent l'article et n'ont pas été reproduits par notre groupe. Partie 2 : 337 mots ; cible inférieure à 4 minutes, à vérifier par répétition chronométrée.
}
\end{frame}

% Slide 11 / 32; intervenant 2

\begin{frame}{Du projet à la réponse du modèle}
\small
\begin{block}{Sélection}
\begin{itemize}
  \item Extraction, filtres et contrôle manuel
\end{itemize}
\end{block}
\begin{block}{Prompt}
\begin{itemize}
  \item Consigne et contexte adaptés à la tâche
\end{itemize}
\end{block}
\begin{block}{Évaluation}
\begin{itemize}
  \item Solutions replacées dans le projet
\end{itemize}
\end{block}
\medskip
{\small Projet complet disponible \(\ne\) projet complet envoyé dans chaque prompt.\par}

\note{
INTERVENANT 2 — Hao QIAN — Contributions et méthode

DIAPOSITIVE 11 / 32

Ils extraient le code et les tests, puis appliquent des filtres et une vérification manuelle. Pour la génération de fonctions, ils demandent notamment une couverture de tests supérieure à cinquante pour cent.

Ensuite, chaque question contient une consigne et les informations utiles à la tâche. Selon le cas, ce sont des imports, des signatures ou un rapport de bug.

Attention : le benchmark fournit le projet complet, mais le prompt ne contient pas forcément tout ce projet. Les solutions générées sont ensuite vérifiées dans leur environnement.

REPÈRES : sélection / contexte / réponse / vérification

SOURCE : Zeng et al. (2024), CoderUJB: An Executable and Unified Java Benchmark for Practical Programming Scenarios. §3.3-3.4, pp. 127-129. DOI: 10.1145/3650212.3652115.

Les chiffres reprennent l'article et n'ont pas été reproduits par notre groupe. Partie 2 : 337 mots ; cible inférieure à 4 minutes, à vérifier par répétition chronométrée.
}
\end{frame}

% Slide 12 / 32; intervenant 2

\begin{frame}{Des critères adaptés à chaque tâche}
\small
\begin{block}{Génération}
\begin{itemize}
  \item 20 propositions par question ; pass-all@1 et pass-all@10
\end{itemize}
\end{block}
\begin{block}{Tests de bugs}
\begin{itemize}
  \item Échec sur la version défectueuse, succès après correction
\end{itemize}
\end{block}
\begin{block}{Détection}
\begin{itemize}
  \item Accuracy : proportion de classifications correctes
\end{itemize}
\end{block}
\medskip
{\small Syntaxe correcte, compilation et réussite des tests sont des étapes différentes.\par}

\note{
INTERVENANT 2 — Hao QIAN — Contributions et méthode

DIAPOSITIVE 12 / 32

Pour les tâches de génération, les auteurs produisent vingt réponses par question. Le score pass-all à k estime la probabilité d’avoir au moins une réponse correcte parmi k propositions. Ils présentent notamment k égal à un et à dix.

Pour une fonction ou une réparation, la réponse doit passer les tests. Pour reproduire un bug, le test doit échouer sur la version défectueuse et réussir sur la version corrigée. La détection utilise le pourcentage de réponses correctes.

Il faut donc comparer des scores de même nature. La personne suivante va présenter les premiers résultats.

REPÈRES : pass-all@k / 20 réponses / accuracy

SOURCE : Zeng et al. (2024), CoderUJB: An Executable and Unified Java Benchmark for Practical Programming Scenarios. §3.5, pp. 129-130. DOI: 10.1145/3650212.3652115.

Les chiffres reprennent l'article et n'ont pas été reproduits par notre groupe. Partie 2 : 337 mots ; cible inférieure à 4 minutes, à vérifier par répétition chronométrée.
}
\end{frame}

\section{Résultats : contexte et modèles}

% Slide 14 / 32; intervenant 3

\begin{frame}{Le contexte aide la génération}
\small
\begin{center}
\small
\renewcommand{\arraystretch}{1.35}
\begin{tabular}{@{}lr@{}}
\toprule
\textbf{Prompt} & \textbf{FCG : pass-all@1} \\
\midrule
Un exemple & 12,84 \% \\
Quatre exemples & 14,50 \% \\
Contexte du projet & 21,91 \% \\
\bottomrule
\end{tabular}
\end{center}
\medskip
{\small CodeLlama-13B. Comparaison de prompts ; deux modèles étudiés dans RQ1.\par}

\note{
INTERVENANT 3 — Jianye SHI — Résultats : contexte et modèles

DIAPOSITIVE 14 / 32

Je vais répondre à deux questions : le contexte aide-t-il le modèle, et comment les modèles se comparent-ils ?

Pour la première question, les auteurs testent deux modèles. Ils comparent un prompt avec du contexte du projet à des prompts contenant un ou quatre exemples.

Pour CodeLlama treize milliards, le score de génération de fonctions atteint environ vingt-deux pour cent avec le contexte. Il est d’environ treize pour cent avec un exemple. Mais ce résultat ne signifie pas qu’il améliore automatiquement toutes les tâches.

REPÈRES : RQ1 / contexte du projet / exemples

SOURCE : Zeng et al. (2024), CoderUJB: An Executable and Unified Java Benchmark for Practical Programming Scenarios. RQ1, tableau 3, p. 131. DOI: 10.1145/3650212.3652115.

Les chiffres reprennent l'article et n'ont pas été reproduits par notre groupe. Partie 3 : 344 mots ; cible inférieure à 4 minutes, à vérifier par répétition chronométrée.
}
\end{frame}

% Slide 15 / 32; intervenant 3

\begin{frame}{Le classement dépend de la tâche}
\small
\begin{center}
\small
\renewcommand{\arraystretch}{1.35}
\begin{tabular}{@{}l@{\hspace{8mm}}r@{\hspace{7mm}}r@{}}
\toprule
\textbf{Modèle} & \textbf{FCG} & \textbf{APR} \\
\midrule
CodeLlama-34B & 22,82 \% & 5,01 \% \\
GPT-3.5-Turbo & 23,37 \% & 9,31 \% \\
GPT-4 & 30,52 \% & 18,76 \% \\
\bottomrule
\end{tabular}
\end{center}
\medskip
{\small Même indicateur : pass-all@1. Versions étudiées : GPT-3.5-Turbo-0301 et GPT-4-0314.\par}

\note{
INTERVENANT 3 — Jianye SHI — Résultats : contexte et modèles

DIAPOSITIVE 15 / 32

Passons à la comparaison des modèles. Ici, les scores utilisent pass-all à un.

En génération de fonctions, CodeLlama trente-quatre milliards est proche de GPT trois point cinq : environ vingt-trois pour cent dans les deux cas. GPT quatre obtient environ trente et un pour cent.

En réparation, l’écart est plus marqué. CodeLlama est à environ cinq pour cent, GPT trois point cinq à neuf, et GPT quatre à dix-neuf.

Ces résultats concernent les versions testées dans l’article. Ils ne constituent pas un classement des modèles actuels.

REPÈRES : RQ2 / mêmes métriques / modèles de l’étude

SOURCE : Zeng et al. (2024), CoderUJB: An Executable and Unified Java Benchmark for Practical Programming Scenarios. Tableaux 2 et 4, pp. 130, 132. DOI: 10.1145/3650212.3652115.

Les chiffres reprennent l'article et n'ont pas été reproduits par notre groupe. Partie 3 : 344 mots ; cible inférieure à 4 minutes, à vérifier par répétition chronométrée.
}
\end{frame}

% Slide 16 / 32; intervenant 3

\begin{frame}{Des difficultés au-delà du code à compléter}
\small
\begin{center}
\small
\renewcommand{\arraystretch}{1.35}
\begin{tabular}{@{}lr@{}}
\toprule
\textbf{GPT-4 : tâche} & \textbf{pass-all@1} \\
\midrule
FCG : compléter une fonction & 30,52 \% \\
CTG : test à partir du code & 24,18 \% \\
ITG : test reproduisant un bug & 6,66 \% \\
APR : réparer une fonction & 18,76 \% \\
\bottomrule
\end{tabular}
\end{center}
\medskip
{\small Les quatre scores utilisent k = 1. Les exigences de réussite varient selon la tâche.\par}

\note{
INTERVENANT 3 — Jianye SHI — Résultats : contexte et modèles

DIAPOSITIVE 16 / 32

Même le meilleur modèle de ce tableau rencontre des difficultés. Pour GPT quatre, le score est d’environ trente et un pour cent pour compléter une fonction. Il descend à vingt-quatre pour cent pour créer un test à partir du code.

Pour produire un test qui reproduit un bug signalé, il atteint seulement six virgule soixante-six pour cent. La réparation atteint environ dix-neuf pour cent.

Les tâches et les critères diffèrent, donc il faut rester prudent dans la comparaison. Mais ces résultats montrent qu’un bon score en génération de fonctions ne résume pas toutes les capacités.

REPÈRES : FCG / CTG / ITG / APR

SOURCE : Zeng et al. (2024), CoderUJB: An Executable and Unified Java Benchmark for Practical Programming Scenarios. Tableau 4, p. 132 ; valeurs choisies plutôt que le texte contradictoire de RQ2. DOI: 10.1145/3650212.3652115.

Les chiffres reprennent l'article et n'ont pas été reproduits par notre groupe. Partie 3 : 344 mots ; cible inférieure à 4 minutes, à vérifier par répétition chronométrée.
}
\end{frame}

% Slide 17 / 32; intervenant 3

\begin{frame}{La détection reste particulièrement difficile}
\small
\begin{center}
\renewcommand{\arraystretch}{1.5}
\begin{tabular}{ccc}
{\Large\structure{940}} & {\Large\structure{50\,\%}} & {\Large\structure{52,16\,\%}} \\
fonctions classées & référence constante & accuracy de GPT-4
\end{tabular}
\end{center}
\medskip
{\small 470 fonctions défectueuses + 470 corrigées. Résultat descriptif, sans preuve de supériorité statistique.\par}

\note{
INTERVENANT 3 — Jianye SHI — Résultats : contexte et modèles

DIAPOSITIVE 17 / 32

La détection est encore un autre problème. Le jeu de données contient autant de fonctions défectueuses que de fonctions corrigées. Une réponse toujours identique donnerait donc cinquante pour cent de bonnes réponses.

GPT quatre obtient cinquante-deux virgule seize pour cent. C’est proche de cette référence. Certains modèles produisent aussi des réponses mal formatées.

Le message principal est donc de regarder les résultats tâche par tâche. Je passe maintenant aux effets de l’entraînement, qui peuvent aussi changer selon la tâche.

REPÈRES : 940 exemples / 50 \% / 52,16 \%

SOURCE : Zeng et al. (2024), CoderUJB: An Executable and Unified Java Benchmark for Practical Programming Scenarios. §3.3, RQ2 et tableau 4, pp. 128, 132. DOI: 10.1145/3650212.3652115.

Les chiffres reprennent l'article et n'ont pas été reproduits par notre groupe. Partie 3 : 344 mots ; cible inférieure à 4 minutes, à vérifier par répétition chronométrée.
}
\end{frame}

\section{Résultats : effets de l’entraînement}

% Slide 19 / 32; intervenant 4

\begin{frame}{Deux façons de poursuivre l’entraînement}
\small
\begin{block}{RQ3 : spécialisation par langage}
\begin{itemize}
  \item Du code supplémentaire en Python ou en Java
\end{itemize}
\end{block}
\begin{block}{RQ4 : ajustement sur instructions}
\begin{itemize}
  \item Des exemples de demandes et de réponses
\end{itemize}
\end{block}
\medskip
{\small Comparaison avec le modèle de départ, selon la tâche et la métrique.\par}

\note{
INTERVENANT 4 — Xiang BIAN — Résultats : effets de l’entraînement

DIAPOSITIVE 19 / 32

Je vais maintenant expliquer les effets de deux types d’entraînement. 

Le premier est le préentraînement supplémentaire dans un langage précis. Par exemple, on continue à entraîner le modèle avec du code Python ou Java.

Le second est l’ajustement sur des instructions. On entraîne alors le modèle à répondre à des demandes avec des exemples de tâches et de réponses.

Ils comparent ces versions aux modèles de départ : quelles tâches progressent ?

REPÈRES : RQ3 / langage / RQ4 / instructions

SOURCE : Zeng et al. (2024), CoderUJB: An Executable and Unified Java Benchmark for Practical Programming Scenarios. §4.1-4.2, RQ3-RQ4, pp. 130, 133. DOI: 10.1145/3650212.3652115.

Les chiffres reprennent l'article et n'ont pas été reproduits par notre groupe. Partie 4 : 349 mots ; cible inférieure à 4 minutes, à vérifier par répétition chronométrée.
}
\end{frame}

% Slide 20 / 32; intervenant 4

\begin{frame}{Le langage d’entraînement change les résultats}
\small
\textbf{StarCoder-15B}
\begin{center}
\small
\renewcommand{\arraystretch}{1.35}
\begin{tabular}{@{}lrr@{}}
\toprule
\textbf{Version} & \textbf{CoderUJB FCG / Java} & \textbf{HumanEval / Python} \\
\midrule
Base & 15,32 \% & 30,35 \% \\
+ Python & 14,39 \% & 33,57 \% \\
+ Java & 18,82 \% & 27,07 \% \\
\bottomrule
\end{tabular}
\end{center}
\medskip
{\small Génération de fonctions, k = 1. Deux benchmarks différents, comparés chacun entre versions.\par}

\note{
INTERVENANT 4 — Xiang BIAN — Résultats : effets de l’entraînement

DIAPOSITIVE 20 / 32

Voici un exemple avec StarCoder. Sur la génération de fonctions Java de CoderUJB, le modèle de base obtient environ quinze pour cent. La version entraînée davantage sur Python descend à environ quatorze pour cent. La version Java monte à environ dix-neuf pour cent.

Sur HumanEval en Python, on observe le mouvement inverse : la version Python progresse et la version Java recule.

Cela illustre un compromis entre langages. Mais pour reproduire des bugs ou réparer des fonctions, les auteurs observent des effets moins réguliers. Il ne faut donc pas généraliser cet exemple à toutes les tâches.

REPÈRES : StarCoder / Python / Java / génération

SOURCE : Zeng et al. (2024), CoderUJB: An Executable and Unified Java Benchmark for Practical Programming Scenarios. RQ3 et tableau 4, pp. 132-133. DOI: 10.1145/3650212.3652115.

Les chiffres reprennent l'article et n'ont pas été reproduits par notre groupe. Partie 4 : 349 mots ; cible inférieure à 4 minutes, à vérifier par répétition chronométrée.
}
\end{frame}

% Slide 21 / 32; intervenant 4

\begin{frame}{Les instructions n’améliorent pas tout}
\small
\textbf{CodeLlama-13B}
\begin{center}
\small
\renewcommand{\arraystretch}{1.35}
\begin{tabular}{@{}lrr@{}}
\toprule
\textbf{Tâche et indicateur} & \textbf{Base} & \textbf{Instruct} \\
\midrule
FCG : pass-all@1 & 21,91 \% & 13,28 \% \\
CTG : pass-all@1 & 12,61 \% & 6,14 \% \\
APR : pass-all@1 & 4,50 \% & 4,09 \% \\
APR : pass-all@10 & 8,39 \% & 15,72 \% \\
\bottomrule
\end{tabular}
\end{center}
\medskip
{\small Un gain à k = 10 ne garantit pas un gain à k = 1.\par}

\note{
INTERVENANT 4 — Xiang BIAN — Résultats : effets de l’entraînement

DIAPOSITIVE 21 / 32

L’ajustement sur instructions donne aussi des résultats contrastés. Regardons CodeLlama treize milliards et sa version Instruct.

Pour compléter une fonction, le score à une tentative passe d’environ vingt-deux à treize pour cent. Pour créer un test à partir du code, il passe d’environ treize à six pour cent.

Pour la réparation, il faut préciser le nombre de tentatives. Avec une seule proposition, le score baisse légèrement. Avec dix propositions, il passe d’environ huit à seize pour cent.

Dire simplement que les instructions améliorent la réparation serait donc trop général. La métrique compte aussi.

REPÈRES : CodeLlama-13B / FCG / CTG / APR

SOURCE : Zeng et al. (2024), CoderUJB: An Executable and Unified Java Benchmark for Practical Programming Scenarios. RQ4 et tableau 4, pp. 132-134. DOI: 10.1145/3650212.3652115.

Les chiffres reprennent l'article et n'ont pas été reproduits par notre groupe. Partie 4 : 349 mots ; cible inférieure à 4 minutes, à vérifier par répétition chronométrée.
}
\end{frame}

% Slide 22 / 32; intervenant 4

\begin{frame}{Une conclusion liée au protocole étudié}
\small
\begin{block}{Constat}
\begin{itemize}
  \item Les effets varient selon la tâche et le nombre d’essais
\end{itemize}
\end{block}
\begin{block}{Interprétation des auteurs}
\begin{itemize}
  \item La proximité avec le préentraînement pourrait compter
\end{itemize}
\end{block}
\begin{block}{Cas particulier}
\begin{itemize}
  \item CodeLlama-Instruct-34B : réponses répétées en FCG et CTG
\end{itemize}
\end{block}
\medskip
{\small Les observations concernent les modèles et les réglages de l’article.\par}

\note{
INTERVENANT 4 — Xiang BIAN — Résultats : effets de l’entraînement

DIAPOSITIVE 22 / 32

Les auteurs proposent une explication : l’effet de l’entraînement pourrait dépendre de la proximité entre la tâche évaluée et la tâche de préentraînement. C’est une interprétation des résultats, pas une règle démontrée pour tous les modèles.

CodeLlama Instruct trente-quatre milliards produit des réponses identiques pour différentes questions dans certaines tâches. Les auteurs parlent d’un effondrement du mode de génération et écartent ces résultats de leur analyse correspondante.

Nous retenons donc qu’il faut tester chaque stratégie sur plusieurs tâches. La personne suivante va préciser les limites de cette étude.

REPÈRES : effets variables / hypothèse / prudence

SOURCE : Zeng et al. (2024), CoderUJB: An Executable and Unified Java Benchmark for Practical Programming Scenarios. RQ3-RQ4, tableau 5 et discussion, pp. 133-134. DOI: 10.1145/3650212.3652115.

Les chiffres reprennent l'article et n'ont pas été reproduits par notre groupe. Partie 4 : 349 mots ; cible inférieure à 4 minutes, à vérifier par répétition chronométrée.
}
\end{frame}

\section{Limites de l’étude}

% Slide 24 / 32; intervenant 5

\begin{frame}{Les trois limites principales}
\small
\begin{block}{Validité interne}
\begin{itemize}
  \item Erreurs possibles dans les scripts
\end{itemize}
\end{block}
\begin{block}{Validité externe}
\begin{itemize}
  \item Choix des modèles et généralisation
\end{itemize}
\end{block}
\begin{block}{Validité des mesures}
\begin{itemize}
  \item Indicateurs et qualité des tests
\end{itemize}
\end{block}
\medskip
{\small Les auteurs cherchent à limiter ces risques.\par}

\note{
INTERVENANT 5 — Zijie Fang — Limites de l’étude

DIAPOSITIVE 24 / 32

Je vais maintenant présenter les limites de l’étude. Je distingue celles que les auteurs discutent explicitement et les réserves que nous ajoutons.

Dans leur section sur la validité, les auteurs présentent trois types de risques. Le premier concerne les erreurs possibles dans les scripts. Le deuxième concerne le choix des modèles et la portée des conclusions. Le troisième concerne les indicateurs utilisés.

Les auteurs cherchent à réduire ces risques, sans pouvoir les supprimer.

REPÈRES : validité interne / externe / des mesures

SOURCE : Zeng et al. (2024), CoderUJB: An Executable and Unified Java Benchmark for Practical Programming Scenarios. §6, p. 134. DOI: 10.1145/3650212.3652115.

Les chiffres reprennent l'article et n'ont pas été reproduits par notre groupe. Partie 5 : 346 mots ; cible inférieure à 4 minutes, à vérifier par répétition chronométrée.
}
\end{frame}

% Slide 25 / 32; intervenant 5

\begin{frame}{Réussir les tests ne suffit pas}
\small
\begin{block}{Ce que font les auteurs}
\begin{itemize}
  \item Vérification des scripts et publication des résultats
\end{itemize}
\end{block}
\begin{block}{Limite des tests}
\begin{itemize}
  \item Un programme peut encore contenir des bugs
\end{itemize}
\end{block}
\begin{block}{Limite de la couverture}
\begin{itemize}
  \item Exécuter du code $\ne$ vérifier tous ses comportements
\end{itemize}
\end{block}
\medskip
{\small Passer les tests $\ne$ correction garantie.\par}

\note{
INTERVENANT 5 — Zijie Fang — Limites de l’étude

DIAPOSITIVE 25 / 32

Pour les erreurs d’implémentation, les auteurs expliquent qu’ils ont relu leurs scripts. Ils publient également le code et les résultats pour permettre des vérifications.

Concernant les mesures, ils utilisent plusieurs indicateurs et vérifient la couverture des tests. C’est utile, mais un programme qui passe les tests peut encore contenir un bug non détecté.

De même, couvrir une ligne signifie qu’un test l’a exécutée. Cela ne garantit pas que le test vérifie correctement tous les comportements. Il faut donc comprendre les scores comme une réussite selon les critères du benchmark, pas comme une preuve de correction totale.

REPÈRES : code public / tests / couverture

SOURCE : Zeng et al. (2024), CoderUJB: An Executable and Unified Java Benchmark for Practical Programming Scenarios. §3.5 et §6, pp. 129-130, 134 ; distinction couverture/correction : analyse du groupe. DOI: 10.1145/3650212.3652115.

Les chiffres reprennent l'article et n'ont pas été reproduits par notre groupe. Partie 5 : 346 mots ; cible inférieure à 4 minutes, à vérifier par répétition chronométrée.
}
\end{frame}

% Slide 26 / 32; intervenant 5

\begin{frame}{Des résultats difficiles à généraliser}
\small
\begin{block}{Modèles}
\begin{itemize}
  \item Des versions précises des LLM
\end{itemize}
\end{block}
\begin{block}{Données}
\begin{itemize}
  \item 17 projets Java issus de Defects4J
\end{itemize}
\end{block}
\begin{block}{Réparation (APR)}
\begin{itemize}
  \item Des bugs dans une seule fonction
\end{itemize}
\end{block}
\medskip
{\small Pas nécessairement les mêmes résultats ailleurs.\par}

\note{
INTERVENANT 5 — Zijie Fang — Limites de l’étude

DIAPOSITIVE 26 / 32

Les auteurs choisissent des modèles qu’ils considèrent comme représentatifs. Mais ce sont des versions précises, évaluées dans les conditions de cette étude.

Nous ajoutons une réserve sur les données. Les questions proviennent de dix-sept projets Java d’une même source. Cela facilite une construction cohérente, mais ne représente pas tous les langages ni tous les logiciels.

De plus, la réparation se concentre sur des défauts dans une seule fonction. On ne peut donc pas supposer que les résultats seraient identiques pour une correction répartie dans plusieurs fichiers.

REPÈRES : Java / 17 projets / versions testées

SOURCE : Zeng et al. (2024), CoderUJB: An Executable and Unified Java Benchmark for Practical Programming Scenarios. §3.2-3.3, §4.2, §6, pp. 127-130, 134 ; portée des projets : analyse du groupe. DOI: 10.1145/3650212.3652115.

Les chiffres reprennent l'article et n'ont pas été reproduits par notre groupe. Partie 5 : 346 mots ; cible inférieure à 4 minutes, à vérifier par répétition chronométrée.
}
\end{frame}

% Slide 27 / 32; intervenant 5

\begin{frame}{Une évaluation encore limitée}
\small
\begin{block}{Protocole étudié}
\begin{itemize}
  \item Prompt $\rightarrow$ réponse $\rightarrow$ évaluation
\end{itemize}
\end{block}
\begin{block}{Autre approche possible}
\begin{itemize}
  \item Tester $\rightarrow$ corriger $\rightarrow$ recommencer
\end{itemize}
\end{block}
\begin{block}{Contexte fourni}
\begin{itemize}
  \item Le modèle ne reçoit pas forcément tout le projet
\end{itemize}
\end{block}
\medskip
{\small Ces limites précisent la portée des conclusions.\par}

\note{
INTERVENANT 5 — Zijie Fang — Limites de l’étude

DIAPOSITIVE 27 / 32

Enfin, les auteurs étudient principalement des réponses à des questions ou des complétions de code. Ils ne mesurent pas toute l’activité d’un développeur assisté par une IA.

Par exemple, une méthode qui exécute ses propres tests puis corrige plusieurs fois sa réponse pourrait se comporter autrement. 

Il faut aussi distinguer le projet disponible dans le benchmark et le contexte réellement fourni au modèle. Une information absente du prompt peut influencer le résultat.

Ces limites précisent donc où les conclusions s’appliquent. La dernière partie va proposer notre évaluation et des améliorations possibles.

REPÈRES : prompt / réponse / outils / portée

SOURCE : Zeng et al. (2024), CoderUJB: An Executable and Unified Java Benchmark for Practical Programming Scenarios. §3.4 et conclusion 2 de RQ2, pp. 128-129, 132 ; portée du protocole : analyse du groupe. DOI: 10.1145/3650212.3652115.

Les chiffres reprennent l'article et n'ont pas été reproduits par notre groupe. Partie 5 : 346 mots ; cible inférieure à 4 minutes, à vérifier par répétition chronométrée.
}
\end{frame}

\section{Critique et conclusion}

% Slide 29 / 32; intervenant 6

\begin{frame}{Notre évaluation : un apport utile}
\small
\begin{block}{Pertinence}
\begin{itemize}
  \item Plusieurs activités réelles de développement
\end{itemize}
\end{block}
\begin{block}{Vérification}
\begin{itemize}
  \item Exécution et critères adaptés aux tâches
\end{itemize}
\end{block}
\begin{block}{Transparence}
\begin{itemize}
  \item Code, scripts et résultats rendus publics
\end{itemize}
\end{block}
\medskip
{\small Notre appréciation : une contribution à la méthode d’évaluation, au-delà d’un classement.\par}

\note{
INTERVENANT 6 — Yupan Huang — Critique et conclusion

DIAPOSITIVE 29 / 32

Pour terminer, voici notre évaluation de l’article. Son principal point fort est de rapprocher l’évaluation de plusieurs activités réelles de développement.

Les auteurs réunissent plusieurs tâches, utilisent des projets existants et vérifient les solutions avec des critères adaptés. Ils rendent aussi leurs données et leurs résultats accessibles. 

Nous apprécions surtout le fait de regarder les résultats par tâche. Un seul score peut masquer des différences importantes. Cet apport concerne aussi la manière d’évaluer les futurs modèles.

REPÈRES : forces / tâches / exécution / transparence

SOURCE : Zeng et al. (2024), CoderUJB: An Executable and Unified Java Benchmark for Practical Programming Scenarios. §3, §5-7, pp. 126-130, 134 ; appréciation du groupe. DOI: 10.1145/3650212.3652115.

Les chiffres reprennent l'article et n'ont pas été reproduits par notre groupe. Partie 6 : 348 mots ; cible inférieure à 4 minutes, à vérifier par répétition chronométrée.
}
\end{frame}

% Slide 30 / 32; intervenant 6

\begin{frame}{Des comparaisons à mieux isoler}
\small
\begin{block}{Constat du protocole}
\begin{itemize}
  \item Base : complétion ; Instruct : dialogue
\end{itemize}
\end{block}
\begin{block}{Réserve}
\begin{itemize}
  \item Le format peut contribuer à l’écart observé
\end{itemize}
\end{block}
\begin{block}{Amélioration proposée}
\begin{itemize}
  \item Tester plusieurs formats et quantifier l’incertitude
\end{itemize}
\end{block}
\medskip
{\small Un contenu équivalent ne garantit pas un effet identique du format.\par}

\note{
INTERVENANT 6 — Yupan Huang — Critique et conclusion

DIAPOSITIVE 30 / 32

Notre première réserve concerne l’interprétation des différences entre modèles. Les modèles de base utilisent un format de complétion. Les modèles ajustés sur instructions utilisent un format de dialogue.

Les auteurs gardent le même contenu et considèrent ces formats comme équivalents. Mais le format peut aussi influencer la réponse. Nous voudrions donc une expérience supplémentaire qui sépare mieux l’effet de l’entraînement et celui du prompt.

Nous proposerions aussi plusieurs répétitions et des intervalles d’incertitude. Ce sont nos propositions pour vérifier la stabilité des écarts.

REPÈRES : format du prompt / entraînement / protocole

SOURCE : Zeng et al. (2024), CoderUJB: An Executable and Unified Java Benchmark for Practical Programming Scenarios. §3.4, §4.2, pp. 128, 130 ; critique et propositions du groupe. DOI: 10.1145/3650212.3652115.

Les chiffres reprennent l'article et n'ont pas été reproduits par notre groupe. Partie 6 : 348 mots ; cible inférieure à 4 minutes, à vérifier par répétition chronométrée.
}
\end{frame}

% Slide 31 / 32; intervenant 6

\begin{frame}{Une lecture attentive des chiffres}
\small
\textbf{GPT-4, tâche ITG}
\begin{center}
\small
\renewcommand{\arraystretch}{1.35}
\begin{tabular}{@{}lr@{}}
\toprule
\textbf{Indicateur} & \textbf{Valeur du tableau 4} \\
\midrule
pass-all@1 & 6,66 \% \\
pass-all@10 & 15,76 \% \\
\bottomrule
\end{tabular}
\end{center}
\medskip
{\small RQ2 mentionne 15,76 parmi des scores annoncés à k = 1. Nos diapositives suivent le tableau 4.\par}

\note{
INTERVENANT 6 — Yupan Huang — Critique et conclusion

DIAPOSITIVE 31 / 32

Autre remarque : certains chiffres sont incohérents. Pour GPT quatre, le tableau quatre donne six virgule soixante-six pour cent en génération de tests à partir d’un bug, avec une tentative.

Le texte du résultat deux mentionne quinze virgule soixante-seize pour cent dans une liste présentée comme des scores à une tentative. Or ce chiffre correspond à dix tentatives dans le tableau.

Nous avons donc utilisé les valeurs du tableau et indiqué k sur nos diapositives. Cette incohérence mérite une correction, mais elle ne suffit pas, à elle seule, à invalider l’ensemble de l’étude.

REPÈRES : tableau 4 / k = 1 / k = 10

SOURCE : Zeng et al. (2024), CoderUJB: An Executable and Unified Java Benchmark for Practical Programming Scenarios. RQ2, p. 131 et tableau 4, p. 132 ; incohérence vérifiée par le groupe. DOI: 10.1145/3650212.3652115.

Les chiffres reprennent l'article et n'ont pas été reproduits par notre groupe. Partie 6 : 348 mots ; cible inférieure à 4 minutes, à vérifier par répétition chronométrée.
}
\end{frame}

% Slide 32 / 32; intervenant 6

\begin{frame}{Conclusion et pistes de prolongement}
\small
\begin{block}{Élargir}
\begin{itemize}
  \item Autres projets, langages et méthodes avec outils
\end{itemize}
\end{block}
\begin{block}{Contrôler}
\begin{itemize}
  \item Risque de contamination et sensibilité au protocole
\end{itemize}
\end{block}
\begin{block}{Retenir}
\begin{itemize}
  \item Comparer par tâche, avec la même métrique
\end{itemize}
\end{block}
\medskip
{\small CoderUJB : une base utile pour une évaluation plus réaliste, dans un périmètre défini.\par}

\note{
INTERVENANT 6 — Yupan Huang — Critique et conclusion

DIAPOSITIVE 32 / 32

Pour prolonger cette étude, nous proposerions des projets plus variés, d’autres langages et une vérification du risque de présence des exemples dans les données d’entraînement. Nous testerions aussi des méthodes capables d’utiliser des outils et de corriger leurs réponses.

En conclusion, CoderUJB montre l’intérêt d’évaluer les modèles dans un contexte de projet et sur plusieurs tâches. Les effets de l’entraînement varient, et les indicateurs doivent être lus avec précision.

Notre bilan est donc positif, avec des réserves sur la portée et sur certaines comparaisons. Merci pour votre attention. Nous sommes prêts à répondre à vos questions.

REPÈRES : plus de diversité / contrôles / évaluation par tâche

SOURCE : Zeng et al. (2024), CoderUJB: An Executable and Unified Java Benchmark for Practical Programming Scenarios. §5-7, p. 134 ; prolongements et bilan du groupe. DOI: 10.1145/3650212.3652115.

Les chiffres reprennent l'article et n'ont pas été reproduits par notre groupe. Partie 6 : 348 mots ; cible inférieure à 4 minutes, à vérifier par répétition chronométrée.
}
\end{frame}

\end{document}
